\PassOptionsToPackage{x11names}{xcolor}
\documentclass{beamer}

\usepackage{agda}
\usepackage{graphicx}
\usepackage{catchfilebetweentags}
\makeatletter

\newrobustcmd*\OrigExecuteMetaData[2][\jobname]{%
\CatchFileBetweenTags\CatchFBT@tok{#1}{#2}%
\global\expandafter\CatchFBT@tok\expandafter{%
\expandafter}\the\CatchFBT@tok
}%\OrigExecuteMetaData

\newrobustcmd*\ChkExecuteMetaData[2][\jobname]{%
\CatchFileBetweenTags\CatchFBT@tok{#1}{#2}%
\edef\mytokens{\detokenize\expandafter{\the\CatchFBT@tok}}
\ifx\mytokens\empty\PackageError{catchfilebetweentags}{the tag #2 is not found\MessageBreak in file #1 \MessageBreak called from \jobname.tex}{use a different tag}\fi%
}%\ChkExecuteMetaData

\renewrobustcmd*\ExecuteMetaData[2][\jobname]{%
\ChkExecuteMetaData[#1]{#2}%
\OrigExecuteMetaData[#1]{#2}%
}

\makeatother

\usepackage[utf8]{inputenc}
\usepackage{newunicodechar}
\usepackage{newtxmath}
%\usepackage{newtxtext}

% Misc symbols
\newunicodechar{⌞}{\ensuremath{\llcorner}}
\newunicodechar{⌟}{\ensuremath{\lrcorner}}
\newunicodechar{′}{\ensuremath{\prime}}
\newunicodechar{−}{\ensuremath{-}}
\newunicodechar{─}{\ensuremath{-}}
\newunicodechar{◆}{\ensuremath{\Diamondblack}}
\newunicodechar{⧫}{\ensuremath{\blacklozenge}}
\newunicodechar{∷}{\ensuremath{::}}
\newunicodechar{∙}{\ensuremath{\bullet}}
\newunicodechar{□}{\ensuremath{\Box}}
\newunicodechar{∎}{\ensuremath{\blacksquare}}
\newunicodechar{⋆}{\ensuremath{\star}}
\newunicodechar{∣}{\ensuremath{|}}

% indices
\newunicodechar{₀}{\ensuremath{_0}}
\newunicodechar{₁}{\ensuremath{_1}}
\newunicodechar{₂}{\ensuremath{_2}}
\newunicodechar{₃}{\ensuremath{_3}}

\newunicodechar{ₑ}{\ensuremath{_e}}
\newunicodechar{ᵢ}{\ensuremath{_i}}
\newunicodechar{ₖ}{\ensuremath{_k}}
\newunicodechar{ₘ}{\ensuremath{_m}}
\newunicodechar{ₙ}{\ensuremath{_n}}
\newunicodechar{ᵣ}{\ensuremath{_r}}
\newunicodechar{ₛ}{\ensuremath{_s}}

\newunicodechar{₊}{\ensuremath{_+}}

% exponents
\newunicodechar{⁺}{\ensuremath{\textsuperscript{+}}}
\newunicodechar{⁻}{\ensuremath{\textsuperscript{-}}}

\newunicodechar{²}{\ensuremath{^2}}

\newunicodechar{ᵈ}{\ensuremath{^d}}
\newunicodechar{ⁱ}{\ensuremath{^i}}
\newunicodechar{ˡ}{\ensuremath{^l}}
\newunicodechar{ʳ}{\ensuremath{^r}}
\newunicodechar{ˢ}{\ensuremath{^s}}

\newunicodechar{ᴬ}{\ensuremath{^A}}
\newunicodechar{ᴮ}{\ensuremath{^B}}
\newunicodechar{ᴵ}{\ensuremath{^I}}
\newunicodechar{ᴿ}{\ensuremath{^R}}
\newunicodechar{ᵀ}{\ensuremath{^T}}
\newunicodechar{ᵁ}{\ensuremath{^U}}
\newunicodechar{ⱽ}{\ensuremath{^V}}

% Dots
\newunicodechar{⋯}{\ensuremath{\cdots}}
\newunicodechar{∶}{\ensuremath{:}}

% Equality symbols
\newunicodechar{∼}{\ensuremath{\sim}}
\newunicodechar{≡}{\ensuremath{\equiv}}
\newunicodechar{≢}{\ensuremath{\not\equiv}}
\newunicodechar{≟}{\mbox{\tiny\ensuremath{\stackrel{?}{=}}}}
\newunicodechar{≈}{\ensuremath{\approx}}

% Ordering symbols
\newunicodechar{≤}{\ensuremath{\le}}
\newunicodechar{≥}{\ensuremath{\ge}}

% Arrows
\newunicodechar{↦}{\ensuremath{\mapsto}}
\newunicodechar{⇑}{\ensuremath{\Uparrow}}
\newunicodechar{→}{\ensuremath{\rightarrow}}
\newunicodechar{←}{\ensuremath{\leftarrow}}
\newunicodechar{⇒}{\ensuremath{\Rightarrow}}
\newunicodechar{⇉}{\ensuremath{\rightrightarrows}}

% Mathematical symbols
\newunicodechar{∂}{\ensuremath{\partial}}
\newunicodechar{∋}{\ensuremath{\ni}}
\newunicodechar{∞}{\ensuremath{\infty}}
\newunicodechar{∀}{\ensuremath{\forall}}
\newunicodechar{∃}{\ensuremath{\exists}}
\newunicodechar{⊢}{\ensuremath{\vdash}}
\newunicodechar{⟨}{\ensuremath{\langle}}
\newunicodechar{⟩}{\ensuremath{\rangle}}
\newunicodechar{⊤}{\ensuremath{\top}}
\newunicodechar{∘}{\ensuremath{\circ}}
\newunicodechar{⊎}{\ensuremath{\uplus}}
\newunicodechar{×}{\ensuremath{\times}}
\newunicodechar{ℕ}{\ensuremath{\mathbb{N}}}
\newunicodechar{⟦}{\ensuremath{\llbracket}}
\newunicodechar{⟧}{\ensuremath{\rrbracket}}
\newunicodechar{∈}{\ensuremath{\in}}
\newunicodechar{↑}{\ensuremath{\uparrow}}
\newunicodechar{¬}{\ensuremath{\neg}}
\newunicodechar{⊥}{\ensuremath{\bot}}
\newunicodechar{↝}{\ensuremath{\leadsto}}
\newunicodechar{↶}{\ensuremath{\curvearrowleft}}
\newunicodechar{↺}{\ensuremath{\circlearrowleft}}
\newunicodechar{⊔}{\ensuremath{\sqcup}}
\newunicodechar{⨆}{\ensuremath{\bigsqcup}}
\newunicodechar{∩}{\ensuremath{\cap}}
\newunicodechar{∪}{\ensuremath{\cup}}


%%%%%%%%%%%%%%%%% LETTERS

% Misc
\newunicodechar{ℓ}{\ensuremath{\ell}}

% Greek uppercase
\newunicodechar{Δ}{\ensuremath{\Delta}}
\newunicodechar{Γ}{\ensuremath{\Gamma}}
\newunicodechar{Σ}{\ensuremath{\Sigma}}
\newunicodechar{Θ}{\ensuremath{\Theta}}
\newunicodechar{Ω}{\ensuremath{\Omega}}

% Greek lowercase
\newunicodechar{α}{\ensuremath{\alpha}}
\newunicodechar{β}{\ensuremath{\beta}}
\newunicodechar{δ}{\ensuremath{\delta}}
\newunicodechar{ε}{\ensuremath{\varepsilon}}
\newunicodechar{φ}{\ensuremath{\phi}}
\newunicodechar{γ}{\ensuremath{\gamma}}
\newunicodechar{ι}{\ensuremath{\iota}}
\newunicodechar{κ}{\ensuremath{\kappa}}
\newunicodechar{λ}{\ensuremath{\lambda}}
\newunicodechar{μ}{\ensuremath{\mu}}
\newunicodechar{ψ}{\ensuremath{\psi}}
\newunicodechar{η}{\ensuremath{\eta}}
\newunicodechar{ρ}{\ensuremath{\rho}}
\newunicodechar{σ}{\ensuremath{\sigma}}
\newunicodechar{τ}{\ensuremath{\tau}}
\newunicodechar{ξ}{\ensuremath{\xi}}
\newunicodechar{ζ}{\ensuremath{\zeta}}
\newunicodechar{Π}{\ensuremath{\Pi}}


% mathcal
\newunicodechar{𝓒}{\ensuremath{\mathcal{C}}}
\newunicodechar{𝓔}{\ensuremath{\mathcal{E}}}
\newunicodechar{𝓕}{\ensuremath{\mathcal{F}}}
\newunicodechar{𝓡}{\ensuremath{\mathcal{R}}}
\newunicodechar{𝓢}{\ensuremath{\mathcal{S}}}
\newunicodechar{𝓣}{\ensuremath{\mathcal{T}}}
\newunicodechar{𝓥}{\ensuremath{\mathcal{V}}}
\newunicodechar{𝓦}{\ensuremath{\mathcal{W}}}


\beamertemplatenavigationsymbolsempty

\title{A Quick Tour of Agda}
\author{Guillaume Allais}
\institute{Hausdorff Research Institute for Mathematics}
\date{May 13$^{th}$ 2024}

\begin{document}

\begin{frame}
    \begin{center}
      \includegraphics[width=.5\textwidth]{agda-logo.png}
    \end{center}
    \maketitle
\end{frame}

\begin{frame}{First, an ad break}

  {\Large Scottish PL and Verification Summer School}

  {\large July 29th to August 2nd}

  \bigskip

  2 core courses, 1 invited lecturer, 2*3 advanced courses
  \begin{itemize}
    \item Category Theory (core, applied, semantics of PL)
    \item Type Theory (core, for implicit complexity)
    \item PL (distributed systems, effects and handlers)
    \item Verification (protocols, syntaxes with binding)
  \end{itemize}

  \bigskip

  \url{https://scottish-pl-institute.github.io/splv/2024-strathclyde/}

\end{frame}

\begin{frame}
    \begin{center}
      \includegraphics[width=.5\textwidth]{agda-logo.png}
    \end{center}
    \maketitle
\end{frame}

\begin{frame}{Agda -- A community}
    \begin{itemize}
      \item<1-> Originated in 1999, designed by Catarina Coquand
      \item<2-> Complete rewrite in 2007 by Ulf Norell
      \item<3-> Developped openly on Github
      \item<4-> Core developpers in Sweden and Belgium
      \item<5-> Two week-long hacking meetings per year
    \end{itemize}
\end{frame}

\begin{frame}{Agda -- A programming language}
    \begin{itemize}
      \item<1-> Pluggable backends (the main ones are all Haskell-based)
      \item<2-> Powerful FFI (the whole Haskell ecosystem)
      \item<3-> Escape hatches
    \end{itemize}
\end{frame}

\begin{frame}{Agda -- A theorem prover}
  \begin{itemize}
    \item<1-> Based on type \textbf{theories} (MLTT, 2 level TT, with \& without K, cubical, strict prop, sized types, TT in colour, ...)
    \item<2-> Foundations à la carte through \texttt{OPTIONS}
    \item<3-> A more conservative \texttt{--safe} mode
  \end{itemize}


  \vfill
  \uncover<4->{
  Some soundness bugs, almost all outside of \texttt{--safe}:
  \url{https://github.com/agda/agda/issues?q=is\%3Aissue+is\%3Aopen+sort\%3Aupdated-desc+label\%3Afalse}}
\end{frame}

\begin{frame}{Small kernel?}
  \begin{center}
    \Huge No!
  \end{center}

\vfill
Well, that's not completely true but it's a hard problem:\\
\url{https://jesper.cx/posts/agda-core.html}\\
\url{https://github.com/jespercockx/agda-core}
\end{frame}

\begin{frame}{Automation?}
  \begin{center}
    \Huge No!
  \end{center}

\vfill
Well, that's not true:
\begin{itemize}
  \item \textbf{Extremely} powerful elaboration machinery
    \begin{itemize}
      \item Dependent pattern matching
      \item Unification to reconstruct missing information
      \item Parametrised modules
    \end{itemize}
  \item A meta-programming mechanism (``typed'' tactics)
  \item A form of dependent type providers
\end{itemize}
\end{frame}

\begin{frame}{Flavours}
  \begin{itemize}
    \item Standard Library (\url{https://github.com/agda/agda-stdlib})
    \item Category Theory (\url{https://github.com/agda/agda-categories})
    \item Cubical Agda (\url{https://github.com/agda/cubical})
    \item A Logical Relation for MLTT (\url{https://github.com/mr-ohman/logrel-mltt})
    \item Type Topology (\url{https://github.com/martinescardo/TypeTopology})
    \item Invoking SMT solvers (\url{https://github.com/wenkokke/schmitty})
  \end{itemize}
\end{frame}

\begin{frame}{First, an anecdote}
  \ExecuteMetaData[stringlytyped]{type}
  \ExecuteMetaData[stringlytyped]{hastype}
  \ExecuteMetaData[stringlytyped]{constructors}
\end{frame}

\begin{frame}{Why?!}
  Agda takes seriously the idea that we should use the
  right data structure to represent the objects of study.

  \begin{itemize}
    \item<1-> Inductive Types
    \item<2-> Inductive Families
    \item<3-> Induction Induction
    \item<4-> Induction Recursion
    \item<5-> Higher Inductive Types
    \item<6-> Higher Inductive Inductive Types
  \end{itemize}
\end{frame}

\end{document}
